\section*{第 \ref{ch:g11:fundamental-interactions} 章——基本相互作用}

\begin{solution}{exo:g11:fundamental-interactions:1}
质子：$10^{-15}$（核子尺度）；水分子：$10^{-10}$
（分子尺度——单个分子处于原子尺度）；沙子：
$10^{-4}$（日常物质）；地球：$10^{7}$（行星）；银河：
$10^{21}$（星系）。单位均为米。
\end{solution}

\begin{solution}{exo:g11:fundamental-interactions:2}
带负电：它\emph{得到}了
$N = \num{4.8e-9}/\num{1.6e-19} = \num{3.0e10}$ 个电子。头发
带 $+\qty{4.8e-9}{C}$ 的电荷，由电荷守恒
（见\cref{prop:g11:fundamental-interactions:conservation}）。
\end{solution}

\begin{solution}{exo:g11:fundamental-interactions:3}
$F = \num{8.99e9} \times \dfrac{\num{2.0e-6} \times \num{3.0e-6}}{0.30^2}
= \dfrac{\num{5.39e-2}}{0.090} \approx \qty{0.60}{N}$，吸引
（异号电荷），两力大小相等、
方向相反。
\end{solution}

\begin{solution}{exo:g11:fundamental-interactions:4}
$F_0/4$；\ $9F_0$；\ $2F_0$；\ $F_0$（分子 $\times 4$，分母
$\times 4$）。
\end{solution}

\begin{solution}{exo:g11:fundamental-interactions:5}
(a) 万有引力；(b) 电磁相互作用；(c) 强相互作用；(d) 弱相互作用；
(e) 电磁相互作用。
\end{solution}

\begin{solution}{exo:g11:fundamental-interactions:6}
$F_E = \num{8.99e9} \times (\num{1.60e-19})^2/(\num{2.0e-15})^2
= \num{2.30e-28}/\num{4.0e-30} \approx \qty{58}{N}$;
$F_G = \num{6.67e-11} \times (\num{1.67e-27})^2/\num{4.0e-30}
\approx \qty{4.7e-35}{N}$。比值 $F_E/F_G \approx \num{1.2e36}$。
强相互作用束缚原子核，引力可忽略。
\end{solution}

\begin{solution}{exo:g11:fundamental-interactions:7}
$F_E = \num{2.30e-28}/(\num{5.3e-11})^2 = \num{2.30e-28}/\num{2.81e-21}
\approx \qty{8.2e-8}{N}$;
$F_G = \num{6.67e-11} \times \num{1.67e-27} \times \num{9.11e-31} /
\num{2.81e-21} \approx \qty{3.6e-47}{N}$。比值 $\approx \num{2.3e39}$：
电磁相互作用把原子束缚在一起。
\end{solution}

\begin{solution}{exo:g11:fundamental-interactions:8}
通过接触，电荷均分：$q_A = q_B = \qty{+4.0e-9}{C}$（总量
仍为 $\qty{8.0e-9}{C}$）。于是
$F = \num{8.99e9} \times (\num{4.0e-9})^2/0.10^2 =
\qty{1.4e-5}{N}$，相互排斥（同号）。
\end{solution}

\begin{solution}{exo:g11:fundamental-interactions:9}
$d = \sqrt{k q^2/F} = \sqrt{\num{8.99e9} \times \num{1.0e-12}/0.90}
\approx \qty{0.10}{m}$。若力缩小为原来的 $9$ 分之一，则距离需为原来的 $3$
倍：$d = \qty{0.30}{m}$。
\end{solution}

\begin{solution}{exo:g11:fundamental-interactions:10}
$F = \num{2.30e-28}/(\num{1.0e-10})^2 = \qty{2.3e-8}{N}$。强
相互作用无贡献：$10^{-10}$~m 是其作用范围的 $10^{5}$ 倍。
分子由电磁相互作用束缚。
\end{solution}

\begin{solution}{exo:g11:fundamental-interactions:11}
$q = \num{e23}/\num{e6} \times \num{1.6e-19} = \qty{1.6e-2}{C}$，每
枚硬币上。$F = \num{8.99e9} \times (\num{1.6e-2})^2/1.0^2 \approx
\qty{2.3e6}{N}$ ——即重量为 $\num{2.3e6}/9.81 \approx
\qty{2.3e5}{kg}$（一列满载的货运列车），仅来自百万分之一的
电子。日常物质必须（而且确实）呈电中性，精度远高于
百万分之一。
\end{solution}

\begin{solution}{exo:g11:fundamental-interactions:12}
$F = \num{2.30e-28}/(\num{1.4e-14})^2 = \num{2.30e-28}/\num{1.96e-28}
\approx \qty{1.2}{N}$ ——在此尺度上仍然极大。不，它们
相距为强相互作用范围的 $14$ 倍；每个核子只与最近的
邻居粘合。质子对共 $\frac{92 \times 91}{2} = 4186$ 个，\emph{全部}
在任意距离上都相互排斥，而粘合力不随尺寸增大。
中子只增加吸引而不增加排斥，并稀释质子；但
排斥力随质子数的平方增长，而粘合力只随
近邻数增长。因此质子数超过约 $80$ 后，任何中子
数都无法使账目平衡——此类原子核会衰变
（见\cref{ch:g11:nucleus-radioactivity}）。
\end{solution}

\begin{solution}{exo:g11:fundamental-interactions:13}
平衡：竖直方向 $T\cos\theta = mg$，水平方向
$T\sin\theta = F$，因此 $F = mg\tan\theta$。此处
$\sin\theta = 3.0/50 = 0.060$，小角度时
$\tan\theta \approx \sin\theta = 0.060$。于是
$F = \num{1.0e-3} \times 9.81 \times 0.060 \approx \qty{5.9e-4}{N}$;
$q = d\sqrt{F/k} = 0.060 \times \sqrt{\num{5.9e-4}/\num{8.99e9}}
\approx \qty{1.54e-8}{C}$; $N = \num{1.54e-8}/\num{1.6e-19} \approx
\num{9.6e10}$ 个元电荷。
\end{solution}

\begin{solution}{exo:g11:fundamental-interactions:14}
$F = \num{6.67e-11} \times \num{5.97e24} \times \num{7.35e22} /
(\num{3.84e8})^2 = \num{2.93e37}/\num{1.47e17} \approx \qty{2.0e20}{N}$.
设 $kQ^2/d^2 = F$：
$Q = d\sqrt{F/k} = \num{3.84e8} \times \sqrt{\num{2.0e20}/\num{8.99e9}}
\approx \qty{5.7e13}{C}$。即
$N = \num{5.7e13}/\num{1.6e-19} \approx \num{3.6e32}$ 个电子，其
质量为 $\num{3.6e32} \times \num{9.11e-31} \approx \qty{3.2e2}{kg}$：
约几百千克的电子，就能起到相当于
$\qty{7.35e22}{kg}$ 的月球那样的作用——电的强大与引力的
微弱，尽在这一个数之中。
\end{solution}

\begin{solution}{exo:g11:fundamental-interactions:15}
$F = \num{2.30e-28}/(\num{2.8e-10})^2 = \num{2.30e-28}/\num{7.84e-20}
\approx \qty{2.9e-9}{N}$；重力
$P = \num{5.9e-26} \times 9.81 \approx \qty{5.8e-25}{N}$；比值
$F/P \approx \num{5e15}$。每个键都有\emph{固定}的强度，但
重力随上方整体质量而增大：岩石堆得越高，底层的载荷
无限增长，而其中的键却保持不变。
约 \qty{e4}{m} 厚的岩石所施加的载荷便会压垮化学键——引力
取胜并非靠单个粒子的强弱，而是靠\emph{叠加}，
即在电荷相消之时。
\end{solution}

\begin{solution}{pb:g11:fundamental-interactions:1}
\textbf{1.} 夸克 $< \qty{e-18}{m}$；核子 $\qty{e-15}{m}$；原子核
$10^{-15}$--$\qty{e-14}{m}$；原子 $\qty{e-10}{m}$；分子
$\qty{e-9}{m}$；日常物体 $\qty{1}{m}$；行星 $\qty{e7}{m}$；
行星系 $\qty{e11}{m}$；星系 $\qty{e21}{m}$。

\textbf{2.} 缩放因子 $10^{5}$：原子放大后
宽 $\qty{1}{cm} \times 10^{5} = \qty{1}{km}$。

\textbf{3.} $(10^{-15}/10^{-10})^3 = 10^{-15}$：物质
有 $\num{99.9999999999999}\,\%$ 为空。

\textbf{4.} $\num{7e-5}/\num{e-10} = \num{7e5}$ 个原子——约
一根头发丝的直径上约有一百万个。

\textbf{5.} 从 $10^{-18}$ 到 $10^{21}$：共 $39$ 个 10 的幂。

\textbf{6.} $F_E = \num{8.99e9} \times (\num{1.60e-19})^2 /
(\num{5.3e-11})^2 \approx \qty{8.2e-8}{N}$.

\textbf{7.} $F_G = \num{6.67e-11} \times \num{1.67e-27} \times
\num{9.11e-31}/(\num{5.3e-11})^2 \approx \qty{3.6e-47}{N}$；比值
$F_E/F_G \approx \num{2.3e39}$。电力维系着原子，引力
则只是旁观者。

\textbf{8.} 若无引力，则无任何可测量的改变（它要小 $10^{39}$
倍）。若无电力，则不存在原子、分子，也就根本没有任何物质。

\textbf{9.} “接触”其实是原子的电子
云之间的电排斥：当原子被拉近至约 $\qty{e-10}{m}$ 时，从来没有任何东西真正接触。
地面支持力与摩擦力，不过是宏观尺度上的库仑定律。

\textbf{10.} $q = \num{e23} \times \num{e-9} \times \num{1.6e-19} =
\qty{1.6e-5}{C}$（每个）；
$F = \num{8.99e9} \times (\num{1.6e-5})^2 \approx \qty{2.3}{N}$——相当于
十亿个电子中仅多出一个所产生的清晰可见的力。日常物质
的电中性远胜于 $10^{-9}$。

\textbf{11.} 电磁相互作用。强相互作用力在
约 $\qty{e-14}{m}$ 以外即超出作用范围——比
单个原子小 $10^{4}$ 倍。

\textbf{12.} $F = \num{8.99e9} \times (\num{1.60e-19})^2 /
(\num{1.0e-15})^2 \approx \qty{2.3e2}{N}$.

\textbf{13.} $a = F/m_p = 230/\num{1.67e-27} \approx
\qty{1.4e29}{m/s^2}$，约为 $\num{1.4e28}$ 倍 $g$：没有任何日常力
能与之相比。

\textbf{14.} 强于 \qty{230}{N}（在 $\qty{e-15}{m}$ 处，原子核
得以存在）；对质子与中子同样作用，与电荷无关（中子
被束缚）；作用范围约为 $\qty{e-15}{m}$（原子核保持很小——这种
胶合力无法延伸得更远）。

\textbf{15.} $F = \num{2.30e-28}/(\num{1.4e-14})^2 \approx
\qty{1.2}{N}$。否：距离为作用范围的 $14$ 倍。斥力作用于
\emph{所有} $\frac{92\times91}{2} = 4186$ 对质子；胶合力
仅结合最近邻，因此体系增大时斥力占优。中子
增加胶合力与间距，却不增加斥力——直到超过约 $80$
个质子后，任何配比都无法平衡，原子核随之衰变。

\textbf{16.} 弱相互作用使中子变为质子（反之亦然），
引起 $\beta$ 衰变——它只负责转化，从不束缚。

\textbf{17.} $F = \num{6.67e-11} \times \num{5.97e24} \times
\num{7.35e22}/(\num{3.84e8})^2 \approx \qty{2.0e20}{N}$.

\textbf{18.} $Q_E = \num{1.8e51} \times \num{1.6e-19} =
\qty{2.9e32}{C}$，$Q_M = \qty{3.5e30}{C}$。两力相等意味着
$f^2\,k\,Q_E Q_M = G M_E M_M$:
\[
f = \sqrt{\frac{\num{2.93e37}}{\num{8.99e9} \times \num{2.9e32}
\times \num{3.5e30}}} = \sqrt{\num{3.2e-36}} \approx \num{1.8e-18} :
\]
若失衡达 \emph{$10^{18}$ 中的两份}，便足以抵消引力。

\textbf{19.} 电荷有正负两种，且会相互抵消（第 18 题
表明抵消何等彻底），因此较强的力在大尺度上自行
失效；质量只有一种符号，会无限累加，且无物可以屏蔽它。

\textbf{20.} 原子核中，强相互作用压过了质子的
斥力（在 $\qty{e-15}{m}$ 处）。从原子到山峦，电磁
相互作用将电子束缚于原子核，使原子彼此结合。行星
及更大尺度：万有引力在电荷相消后无所抗衡，维系着行星、
轨道与星系。世界之所以不坍塌，是因为每一层
都有比载荷更强的键支撑，并且不会四散飞开，
因为物质的电中性精确到约 $10^{-18}$——所以每一尺度
恰有一种力占主导，并在必要之处表现为吸引。
\end{solution}
